\subsection{A direct Jacobian estimate}
\label{subsec:deg-raw-jacobian}

In the following two sign ranges, the graph Jacobian itself controls the
complete derivative. An exact logarithmic derivative separates the
endpoint factor from an interior term whose integral has a uniform upper
bound, so no action compactification is needed.

\begin{lemma}\label{deg:lem-raw-jacobian}
Assume either
\(\alpha _0\le0,\alpha _1\le0\), or
\(\alpha _0=-s^2<0,\alpha _1=sp>0,\gamma>0\) and
\[
 {s\over\gamma^{3/2}}\le k_0,\qquad
 {p\over\sqrt\gamma}\le\pi_0<\sqrt {12}.
\]
On the physical input interval selected by Lemma~\ref{lem:loc-capture},
the central transition has a uniformly bounded positive derivative.
After the final choice of the finite saddle--node multiplier, the
complete return is a contraction.
\end{lemma}

\begin{proof}
Put \(\xi=v-\gamma\), the local spatial coordinate of
\eqref{deg:central-field}, and set
\[
 M=-P(\gamma+\xi)
 =-\alpha _0-\alpha _1\xi+3\gamma \xi^2+\xi^3,\qquad
 \mathscr D=M+\ell z=-\dot v.
\]
Along a through orbit, regarded as a graph \(z=z(v)\), write
\[
 f={z\{q_\gamma \xi+\xi^2+Cvz\}\over\mathscr D},\qquad
 \nu={6\gamma\over q_\gamma},\qquad
 N=-\alpha _1+(3-\nu)\xi^2.
\]
Direct differentiation gives
\begin{equation}
 f_z={d\over dv}\log {z^{1+\nu}\over\mathscr D}
       +{N+\{\ell'(v)+(1-\nu)Cv\}z\over\mathscr D}.
 \label{deg:raw-log-jacobian}
\end{equation}
If \(\alpha _1\le0\), then \(N\ge|\alpha _1|+2\xi^2\) after shrinking
\(|\gamma|\).  Let \(\kappa=\frac12\inf\ell\).  On
\(\mathscr D\ge\kappa z\), the last quotient in
\eqref{deg:raw-log-jacobian} is bounded below.  On the complement,
\(M<0\),
\[
 z<{|M|\over\ell-\kappa},\qquad
 |M|\le C(\delta+|\gamma|)N.
\]
Choosing the fixed raw width \(\delta\) and then the coefficient
neighbourhood makes the numerator at least \(N/2\) there.  Thus its
integral in the physical direction from the positive to the negative
section is bounded above.

In the second case put \(\xi=(s/\sqrt\gamma)\widehat x\), with
\(\pi=p/\sqrt\gamma\) and \(k=s/\gamma^{3/2}\).  Then
\[
 M=s^2\{1-\pi\widehat x+3\widehat x^2+k\widehat x^3\}.
\]
Where \(N<0\), the quadratic identity
\[
 1-\pi\widehat x+3\widehat x^2
 =3\left(\widehat x-{\pi\over6}\right)^2+1-{\pi^2\over12}
\]
gives a positive multiple of \(1+\widehat x^2\).  On \(N<0\) one has
\(\widehat x^2\le C\pi/k\), hence
\(|k\widehat x|\le C\sqrt{\pi_0k_0}\).
The cubic term is therefore absorbed by decreasing \(k_0\).
Since \(-N\le sp\), the change of variable cancels the prefactor
\(sp\,s/(s^2\sqrt\gamma)=\pi\le\pi_0\), and hence
\[
 \int_{\{N<0\}}{-N\over\mathscr D}\,d\xi
 \le C\int_{\mathbb R}{d\widehat x\over1+\widehat x^2}<\infty .
\]
Where \(N\ge0\), the preceding two-denominator argument applies.  Indeed
\(M<0\) can then occur only at \(\xi=-u\), \(u>3\gamma\), and
\(N\ge cu^2\), \(|M|\le u^3\le\delta u^2\).

Integration of \eqref{deg:raw-log-jacobian} yields
\[
 0<T_\theta'\le
 C{(z^{1+\nu}/\mathscr D)_-\over
       (z^{1+\nu}/\mathscr D)_+}\le C_0.
\]
The two endpoints lie in the fixed positive-height windows, which bound
the quotient in this estimate.
The estimate holds on the full interval
\(\mathcal T_\theta\cap\operatorname{int}K_+^\sharp
\cap T_\theta^{-1}(\operatorname{int}K_-^\sharp)\), including all
intervening captured inputs. The bounded divided closing of
Lemma~\ref{phys:lem-c3-adapters}, followed by a final multiplier
shrink, gives the contraction.
\end{proof}
